\documentclass{ctexrep}

\usepackage{amsmath,amsthm,extarrows,color,enumitem,framed,bbold,amscd}
\usepackage[a4paper, left = 2cm, right = 2cm]{geometry}
\usepackage{mathpartir}

\usepackage{tikz-cd}

\usepackage{ulem}

\usepackage{enumitem}

\usepackage{hyperref}
\hypersetup{
     colorlinks   = true,
     linkcolor=blue,
     citecolor=red,
     %filecolor=magenta,      % color of file links
     urlcolor=cyan
}
\usepackage{cleveref}
\usepackage{microtype}

\usepackage{graphicx}
\graphicspath{{./images/}}

\theoremstyle{plain}
\newtheorem{theorem}{定理}[section]
\newtheorem{theorem*}{定理}
\newtheorem{proposition*}[theorem*]{命题}
\newtheorem{corollary*}[theorem*]{推论}
\newtheorem{lemma*}[theorem*]{引理}
\newtheorem{definition*}[theorem*]{定义}
\newtheorem{corollary}[theorem]{推论}
\newtheorem{lemma}[theorem]{引理}
\newtheorem{proposition}[theorem]{Proposition}

\theoremstyle{definition}
\newtheorem{definition}[theorem]{定义}
\newtheorem{example}[theorem]{例}

\theoremstyle{remark}
\newtheorem{remark}[theorem]{评论}
\newtheorem*{perspective}{Perspective}
%\newtheorem{definition*}{定义}
%\newtheorem*{remark}{注意}

\providecommand{\tightlist}{%
  \setlength{\itemsep}{0pt}\setlength{\parskip}{0pt}}

\newcommand{\bfP}{\mathbf{P}}
\newcommand{\calC}{\mathcal{C}}
\newcommand{\calE}{\mathcal{E}}
\newcommand{\calS}{\mathcal{S}}
\newcommand\calF{\mathcal{F}}
\newcommand{\yoneda}{\mathbf{y}}

\newcommand{\VarTop}{\mathbf{VarTop}}
\newcommand{\Comm}{\mathbf{Comm}}
\newcommand{\id}{\mathrm{id}}
\newcommand{\Aff}{\mathbf{Aff}}
\newcommand{\Diff}{\mathbf{Diff}}
\newcommand{\Open}{\mathbf{Open}}
\newcommand{\op}{\mathrm{op}}
\newcommand{\true}{\mathrm{true}}
\newcommand{\UG}{\mathbf{U}}
\newcommand{\rmH}{H}
\newcommand{\Gposet}{G\text{-}\mathbf{poset}}
\newcommand{\calP}{\mathcal{P}}
\newcommand{\PG}{\calP\rtimes G}
\newcommand{\kPG}{k\PG}
\newcommand{\EI}{\mathrm{EI}}
\newcommand\frakm{\mathfrak{m}}
\newcommand\pgrhp{\pi_{\gamma}r_P^H}
\newcommand\barN{\overline{N}}
\newcommand\Cl{\mathrm{Cl}}
\newcommand\Tr{\mathrm{Tr}}
\newcommand{\abs}[1]{\left\lvert\#1\right\rvert}
\newcommand\calO{\mathcal{O}}
\newcommand\calU{\mathcal{U}}
\newcommand\OGb{\OG b}
\newcommand\ZOG{\mathbf{Z}\OG}
\newcommand\OGG{(\OG)^G}
\newcommand\Irr{\mathrm{Irr}}
\newcommand\Proj{\mathrm{Proj}}
\newcommand\ksng{k_{\sharp}\widehat{\overline{N}}_G}
\newcommand\LP{\mathcal{LP}}
\newcommand\Res{\mathrm{Res}}
\newcommand\ResNCGPVg{\Res_{\overline{C}_G(P)}^{\overline{N}_G(P_{\gamma})}(V(\gamma))}
\newcommand\BOG{\mathbf{B}(\OG)}
\newcommand\pSubGrp{\mathrm{psGrp}}
\newcommand\GL{\mathrm{GL}}
\newcommand\PGL{\mathrm{PGL}}
\newcommand\Inn{\mathrm{Inn}}
\newcommand\std{\mathrm{std}}
\newcommand\bfZ{\mathbf{Z}}
\newcommand\kP{k\calP}
%\newcommand\Mor{\mathrm{Mor}}
\newcommand\Bas{\mathrm{Bas}}
\newcommand\Orb{\mathrm{Orb}}
\newcommand\calB{\mathcal{B}}
\newcommand\Br{\mathrm{Br}}
\newcommand\Image{\mathrm{Im}}
\newcommand\calD{\mathcal{D}}
\newcommand\Ker{\mathrm{Ker}}
\newcommand\bbF{\mathbb{F}}
\newcommand\bbI{\mathbb{I}}
\newcommand\ttI{\mathtt{I}}
\newcommand\Rad{\mathrm{Rad}}
\newcommand\Soc{\mathrm{Soc}}
\newcommand\Perm{\mathrm{Perm}}
\newcommand\bbA{\mathbb{A}}
\newcommand\rmP{\mathrm{P}}
\newcommand\Fam{\mathbf{Fam}}
\newcommand\rmOb{\mathrm{Ob}}

\newcommand\rmMor{\mathrm{Mor}}
\newcommand\rmType{\mathrm{Type}}
\newcommand\rmTerms{\mathrm{Terms}}
\newcommand\Ty{\mathtt{Ty}}
\newcommand\Ter{\mathtt{Ter}}
\newcommand\sfp{\mathsf{p}}
\newcommand\sfq{\mathsf{q}}
\newcommand\bfone{\mathbf{1}}
\newcommand\Psh{\mathrm{Psh}}
\newcommand\Set{\mathbf{Set}}
\newcommand\C{\mathbf{C}}
\newcommand\Sets{\mathbf{Sets}}
\newcommand\Cube{\square^{\op}}
\newcommand\hatcube{\widehat{□}}
\newcommand\hatsquare{\widehat{□}}
\newcommand\cubicalsets{\widehat{□}}
\newcommand\cSets{[\square^{\op},\Sets]}
\newcommand\Setspointed{\mathbf{Sets}_{\kappa}^{\bullet}}
\newcommand{\qq}{\mathtt{q}}
\newcommand{\pp}{\mathtt{p}}
\newcommand{\Cof}{\mathtt{Cof}}
\newcommand{\dec}{\mathrm{dec}}
\newcommand{\fib}{\mathtt{fib}}
\newcommand{\ttzero}{\mathtt{0}}
\newcommand{\ttone}{\mathtt{1}}

\DeclareMathOperator\Obj{Obj}
\DeclareMathOperator\inl{inl}
\DeclareMathOperator\inr{inr}
\DeclareMathOperator\ob{ob}
\DeclareMathOperator\pt{pt}
\DeclareMathOperator\osf{sf}
\DeclareMathOperator\supp{supp}
\DeclareMathOperator\Sub{Sub}
\DeclareMathOperator\Mor{Mor}
\DeclareMathOperator\Colim{Colim}
\DeclareMathOperator\colim{colim}
\DeclareMathOperator\im{im}
\DeclareMathOperator\Pre{Pr}
\DeclareMathOperator\Hom{Hom}
\DeclareMathOperator\Sh{Sh}
\DeclareMathOperator\Spec{Spec}
\DeclareMathOperator\End{End}
\DeclareMathOperator\Aut{Aut}
\DeclareMathOperator\Func{Func}
\DeclareMathOperator\cof{cof}
\DeclareMathOperator\obj{obj}

\newcommand{\DM}{\mathbf{DM}}
\newcommand{\dM}{\mathsf{dM}}
\newcommand{\yc}{\yoneda c}
\newcommand{\yb}{\yoneda b}
\newcommand{\yd}{\yoneda d}
\newcommand{\yone}{\yoneda[1]}
\newcommand{\one}{\mathbf{1}}
\newcommand{\bbmone}{\mathbbm{1}}
\newcommand{\two}{\mathbbm{2}}

\newcommand{\ax}[1]{\mathbf{ax_{#1}}}
\newcommand{\ul}{{\rule[0ex]{0.5em}{0.4pt}}} % [0.2ex]：基线以上的高度（可调）{1em}：长度（可调）{0.4pt}：粗细
%\newcommand{\ul}{{\scalebox{0.4}[0.4]{\ul }}}
\newcommand{\mi}{{\rule[0.5ex]{0.4em}{0.4pt}}}
\newcommand{\II}{\mathbb{I}}
\newcommand{\Om}{\Omega}
\newcommand{\cube}{\square}
\newcommand{\tti}{\mathtt{i}}

\newcommand{\forces}{\Vdash} % Kripke-Joyal forcing symbol
\newcommand{\valid}{\vdash} % Internal validity symbol

\newcommand{\coloneqq}{\mathrel{:=}}
\newcommand{\Type}{\mathsf{type}}

%\setlength{\parindent}{0pt}

\title{有向类型论中的纤维化}
\author{Benno Lossin}
\date{}

\begin{document}
\maketitle

\begin{abstract}
我们在 Riehl 和 Shulman 发展的综合\textbf{单纯类型论}中研究 \(\infty\)-范畴。具体而言，我们定义了\textbf{余笛卡尔纤维化}，并利用 \textbf{LARI 伴随}与\textbf{始截面}之间的一个新等价性，证明了它们的封闭性质。我们使用实验性证明辅助工具 \textsf{rzk} 对工作进行形式化，并将成果整合到 Riehl 等人的形式化工作中。除新工作外，本文还介绍了一般类型论、同伦类型论以及论文其余部分所使用的单纯类型论。
\end{abstract}

\section*{致谢}

我要感谢我的导师 Tashi Walde，他就本论文的内容进行了大量讨论并付出了大量时间。他帮助我理解和深入这一理论，使我不必担心迷失方向。我还要感谢我的导师 Claudia Scheimbauer，感谢她将我介绍给 Tashi，并指导本论文的完成。

\tableofcontents
\clearpage

\chapter{引言}

现代类型论将几个特殊的概念和思想汇集到一个基础数学理论中，取代了逻辑学和集合论。它与数学中许多基本且非常重要的学科相关联：高阶范畴论、同伦论、序理论、逻辑学、基础数学、形式化数学和构造性数学。它植根于基础数学、逻辑学和理论计算机科学，并通过其与形式化的良好兼容性提供了高度的正确性。此外，它还提供了最直观的数学诠释之一。与同伦论和高阶范畴的联系直到最近才被发现，这导致了对同伦类型论（HoTT）这一新分支的许多有趣探索，最终形成了无公理基础计划的工作[1]。

本论文总结了由 E. Riehl 和 M. Shulman [2] 发展的单纯类型论，作为同伦类型论的扩展。E. Riehl 等人[3] 使用由 N. Kudasov、A. Abounegm 和 D. Danko [4] 基于 Riehl 和 Shulman 理论创建的证明辅助工具 \textsf{rzk} 将该理论形式化。U. Buchholtz 和 J. Weinberger [5] 以及 C. Bardomiano Martinez [6] 对该理论的若干扩展也已被 E. Riehl 等人[3] 部分形式化。我们通过在本论文中形式化余笛卡尔纤维化并将其贡献到 [3] 来继续他们的工作。我们特别遵循了 U. Buchholtz 和 J. Weinberger [5] 以及 J. Weinberger [7] 的工作，他们将余笛卡尔纤维化作为 LARI 族来研究。我们通过证明始截面与 LARI 伴随是等价的来将这些概念联系起来，这简化了形式化工作。我们还证明了余笛卡尔纤维化在乘积、拉回和复合下是封闭的。我们的形式化工作可以在线访问：https://github.com/BennoLossin/sHoTT，当它们被上游接受后，也将在 https://rzk-lang.github.io/sHoTT 以更易读的形式提供。

在本节的剩余部分，我们简要概述类型论中最重要的历史发展。然后我们讨论与其他数学分支的联系，并探讨综合方法的优点。我们还讨论了在同伦类型论中添加方向性的两种不同方式，以便研究高阶范畴。最后，我们概述本论文的结构。

\section{历史}

类型论起源于基础数学以及1910年 A. N. Whitehead 和 B. Russel [8] 在《数学原理》中对罗素悖论的解决。后来，A. Church [9] 在1940年发展了 $\lambda$ 演算，作为一个“数学图灵机”，并引入类型以确保终止。此时，类型的概念已经具有与今天非常相似的含义。理论的另一个巨大飞跃是由 P. Martin-Löf 等人[10] 在1984年通过引入直觉主义类型论实现的。他的开创性工作至今仍是许多类型论的基础，通常被称为 Martin-Löf 类型论（MLTT）。该理论倾向于构造性数学，并成为一个适当的基础系统。MLTT 被许多逻辑学家、数学家和理论计算机科学家进一步发展。我们关注的下一个里程碑是与同伦论的联系以及同伦类型论这一新分支的形成，该分支允许研究综合的 $\infty$-群胚。V. Voevodsky [11] 在2010年提出了幺半等价性的思想，允许将弱等价的对象视为相等。这最终在2013年由无公理基础计划[1] 出版了“HoTT 书”。该书全面概述了同伦类型论，并在其中推导出若干数学结果。该理论的许多进一步扩展已被探索。特别是，E. Riehl 和 M. Shulman [2] 在2017年通过一个严格的定向区间类型添加了方向性。这允许该理论以综合的方式定义 $(\infty,1)$-范畴。

我们的工作核心基于两篇论文：U. Buchholtz 和 J. Weinberger [5] (2021) 以及 J. Weinberger [7] (2024)。他们在 Riehl 和 Shulman 的理论中从几个不同方向考虑了余笛卡尔纤维化的概念。

\section{与其他学科的联系}

数学是一种写下正确内容的句法游戏。理论的规则告诉我们哪些句法构造是有效的。类型论不仅描述了哪些句法构造（也称为项）是有效的，而且为每个项关联了一个类型。类型也是一种句法构造，但处于与项不同的层次。这个类型描述了该项所代表的值的“能力”。例如，函数 $f:A\to B$ 的类型告诉我们哪些值可以作为输入，以及可能得到哪些输出值。

数学家经常谈论值的“能力”，例如“这个函数是弱等价”或“这个对象在 XYZ 范畴中是始对象”。在类型论中，这些陈述无非就是该值（或一个派生值）具有某种类型。当数学家论证时，他们经常回忆并利用这些性质，例如“因为这个函数是等价”或“因为这个对象在 XYZ 中是始对象”。这对应于按照类型所规定的方式使用该值。

类型论自然地形式化了这种“能力”的概念，并使使用它们成为做数学的主要方式。类型论的核心信条是：

\begin{quotation}“类型是命题，项是证明。”\end{quotation}

这个主题是连接大多数其他学科与类型论的纽带。为命题给出证明归结为写出正确类型的项。跟踪不同类型的不同元素的能力是连接类型论与构造性数学和证明相关数学的纽带。一个类型的元素先验地可以是不同的，允许同一陈述有多个不同的证明。这给了我们与同伦论和高阶范畴的联系，因为我们可以使用项跟踪额外的结构，并且它们的非平凡相等性自然地出现在理论中。

\subsection{基础数学与逻辑学}

与之前许多基础系统相比，类型论是特殊的，因为它不是在逻辑学中表述的。类型论本身自然地包含了逻辑构造，并且直接使用逻辑学家定义逻辑的语言来表述。Curry-Howard 元理论同构描述了逻辑学和类型论之间的关系。表1简要概述了类型论如何整合逻辑学。我们将在第2节更详细地讨论这一点。与集合论相比，类型论更适合作为高阶范畴的基础系统。这是因为一个累积的、越来越大的宇宙层次是内置的，并且相等不是一个简单的“$x$ 等于 $y$”的概念，而是包含了两个元素如何相等。此外，集合论没有内置的机制来跟踪高阶结构，而类型论则有。我们注意到，类型论不是一个单一的理论，而是一组可以用来构建理论的原则。就像数学中没有单一的逻辑学，而是有不同的变体一样。

\begin{table}[htbp]
\centering
\renewcommand{\arraystretch}{1.3} % 增加行高以贴合图片排版
\begin{tabular}{|c|c|}
\hline
\textbf{逻辑概念} & \textbf{类型论概念} \\
\hline
命题 & 类型 \\
\hline
命题 $P$ 的证明 & 类型 $P$ 的项 \\
\hline
真 $\top$ & 单元类型 $\mathbb{1}$ \\
\hline
假 $\bot$ & 空类型 $\mathbb{0}$ \\
\hline
合取 $A \land B$ & 乘积类型 $A \times B$ \\
\hline
析取 $A \lor B$ & 不交并类型 $A + B$ \\
\hline
蕴涵 $A \to B$ & 函数类型 $A \to B$ \\
\hline
存在量词 $\exists x.P$ & $\Sigma$-类型 $\Sigma(x : X, P)$ \\
\hline
全称量词 $\forall x.P$ & $\Pi$-类型 $(x : X) \to P$ \\
\hline
\end{tabular}
\caption{逻辑与类型理论之间的 Curry-Howard 元理论同构}
\end{table}


\subsection{形式化数学}

类型论以非常特定的方式使用了逻辑学家的语言。类型论的规则以这样一种方式书写：每条规则都可以机械地应用，而不需要任何数学思想或选择。所谓的类型检查，规则仅决定哪些项属于哪个类型，哪些项无效。这使得它们非常适合推理和开发能够自动检查陈述正确性的程序。形式化数学增加了我们对陈述正确的信任度。此外，类型论方法将形式化数学置于中心位置，而在经典数学中，形式化往往是被事后考虑的。

大多数证明辅助工具建立在 P. Martin-Löf 等人[10] 发展的 Martin-Löf 类型论之上，我们在第1.1节已经提到过。证明辅助工具有例如 Lean (L. de Moura, S. Kong, F. Doorn, and J. Raumer [12], L. d. Moura and S. Ullrich [13] and Lean Contributors [14])、Rocq (T. Coquand and G. Huet [15] and The Coq Development Team [16]) 以及 Isabelle/HOL (T. Nipkow and M. Wenzel [17] and Isabelle Contributors [18])。

对于 E. Riehl 和 M. Shulman [2] 的单纯类型论，有一个实验性的证明辅助工具，称为 \textsf{rzk}，由 N. Kudasov、A. Abouneggm 和 D. Danko [4] 创建。我们使用它来形式化本论文的大部分内容，其基础是 E. Riehl 等人[3] 的工作。

证明辅助工具是发展良好的高阶范畴理论的非常有用的工具。这是因为高阶范畴是一个非常复杂的主题，在处理它们时通常需要记住大量数据。此外，在某些情况下，它们可能比普通范畴更不直观。证明辅助工具有助于避免这些陷阱并简化流程。

\subsection{计算机科学}

毫无疑问，类型论与计算机科学有着最强的联系之一。编程语言的研究——使用各种类型论——是计算机科学的一个完整学科，我们无法在这短短的一段中充分描述它。函数式编程语言基于 Church 的 $\lambda$ 演算，许多现代语言都借鉴了它的概念。也有许多使用深奥类型论的实验性编程语言。计算和证明陈述通过 Curry-Howard 同构密不可分地联系在一起。

\subsection{构造性数学与证明相关性}

构造性数学关注于使用受限的基础系统，该系统不允许凭空产生值。例如，它们不假设选择公理。更令人惊讶的是，它们也不假设排中律，即 $A \lor \neg A$。在经典数学中，这是反证法中常用的不可或缺的工具。从构造性的角度拒绝这个公理的原因是，从公理 $A \lor \neg A$ 无法推导出两者中哪一个为真。这导致反证法是非构造性的。如果通过反证法证明了一个对象 $x$ 的存在性，那么是通过证明它的不存在性会导致矛盾来实现的，从而从未真正构造出这个 $x$。

\begin{example}
展示这种行为的一个简单例子是以下陈述：
\begin{quotation}
是否存在两个无理数 $x$ 和 $y$ 使得 $x^{y}$ 是有理数？
\end{quotation}
答案是肯定的，我们可以使用排中律轻松证明。我们知道 $\sqrt{2}$ 是无理数。现在考虑 $x = y = \sqrt{2}$ 以及下面的论证。要么 $\sqrt{2}^{\sqrt{2}}$ 是无理数，要么是有理数。如果它是有理数，我们就完成了。否则，考虑 $x = \sqrt{2}^{\sqrt{2}}$ 和 $y = \sqrt{2}$，那么它们满足该陈述，因为：

\[
x^{y} = \left(\sqrt{2}^{\sqrt{2}}\right)^{\sqrt{2}} = \sqrt{2}^{2} = 2
\]

从这个论证中，我们不清楚两种情况中哪一种成立，只知道其中一种成立。因此，构造主义者拒绝这种论证。
\end{example}

除了采纳构造性立场的哲学原因外，实际原因之一是任何证明都可以转化为产生定理结果的算法。这在计算机科学领域以及将数学应用于现实世界时特别有用。

类型论通过“类型作为命题”的原则邀请构造主义。我们通过必须创建的项来记录一个命题是如何被证明的，以展示一个类型的居留性。$\Sigma$-类型（类型论中对应存在量词 $\exists x.P(x)$ 的概念）通过构造该 $\Sigma$-类型的项来证明其居留性。这种构造需要两个元素：第一个是见证者 $x$，第二个是见证者满足条件 $P(x)$ 的证明。因此，如果给定一个存在量词的证明，我们总是可以获得一个具体的满足条件的元素。类似地，对于合取，我们需要提供两个输入陈述的证明元素来证明它们的合取。

这使得我们的理论不仅是构造性的，而且是证明相关的。我们关心一个命题可以被证明的方式。可能存在多种不同的方式。这也是同伦论的核心，它研究两个事物如何相等，以及如果它们相等，它们的相等性之间如何相互关联。

我们注意到，我们的理论仍然允许经典数学发生，通过只考虑我们称为命题的类型。这些类型最多有一个元素，这仅对这些类型恢复证明无关性。我们后面会注意到，我们允许对命题假定排中律。

\subsection{同伦论}

同伦论起源于拓扑学，并与高阶范畴论有联系。它的主要研究对象是同伦的概念。同伦本质上是两个拓扑映射 $f,g:X\to Y$ 之间的变形。经典地，它是一个连续映射 $H:X\times [0,1]\to Y$，满足 $H(- ,0) = f$ 和 $H(- ,1) = g$。在这种情况下，人们写作 $f\simeq g$ 并说它们是同伦的。这是一个“等价关系”，但是同一对函数的两个同伦 $H_{1},H_{2}:f\simeq g$ 不一定是相同的。事实上，可以研究同伦之间的同伦 $\alpha :X\times [0,1]\times [0,1]\to Y$，满足 $\alpha (- , - ,0) = H_{1}$ 和 $\alpha (- , - ,1) = H_{2}$。重复这个过程揭示了 $\infty$-群胚的无限结构。

由无公理基础计划[1] 引入的同伦类型论非常好地描述了这种情况。它是拓扑学中分析同伦论的综合复制品。该理论的规则可以写在一张 A5 纸上，更高的结构优美地从相等类型的消去规则中推导出来。我们后面会看到，由于这种更高结构，一个只有一个被证明元素的类型可能非常有趣。这个示例类型类似于 $S^{1}$，并包含预期的更高结构群 $\mathbb{Z}$。

\subsection{序理论、范畴论与高阶范畴论}

在 R. Harper 等人[19] 的讲座中，建立了序理论、普通范畴论和类型论之间的有趣联系。许多概念在这三个理论中都能找到，类似于我们之前看到的 Curry-Howard 元理论同构。例如，乘积类型对应于范畴论中的乘积，在序理论中这对应于交运算。这对偶和 Yoneda 引理也是如此。此外，Heyting 代数自然地从基本类型论中产生，是序理论的研究对象之一。它们也自然地作为特殊范畴存在，并出现在证明论中。

这项研究的自然延续是将其从 $\infty$-群胚扩展到 $\infty$-范畴。我们现在讨论两种可能的方法来添加方向性并使态射不可逆。

\section{向 HoTT 添加方向性的两种方法}

我们已经提到，HoTT 以特别简单而优美的方式表示了同伦论。对于添加方向性并在类型论中展示综合 $\infty$-范畴的方法，情况则不同。虽然这两种理论本身很有用，但与 HoTT 和 $\infty$-群胚的情况相比，它们有所欠缺。产生类似的情况将是该研究领域的“圣杯”，并且很可能不存在。在本节中，我们介绍两种方法，它们各自都有缺陷。第一种方法违反了某些范畴论的预期，第二种方法不建立在 HoTT 之上，而是需要修改，使其本质上不同于 Martin-Löf 类型论。

在本论文中，我们研究第一种方法，并在第3.8节详细解释违反了哪个范畴论预期。这种方法的好处是 HoTT 中的所有陈述仍然成立。此外，使用证明辅助工具 \textsf{rzk} [4] 的形式化工作涵盖了[3] 中相当一部分理论，包括 U. Buchholtz 和 J. Weinberger [5] 以及 C. Bardomiano Martinez [6] 的工作。

\subsection{用形状扩展 HoTT}

我们非常简要地总结 E. Riehl 和 M. Shulman [2] 添加方向性的方法：我们在类型论中添加另一层包含形状的内容，本质上是单纯形 $\Delta^n$ 和某些子形状，例如它们的边界。然后我们将 $\hom_{A}(x,y)$ 定义为从 $\Delta^{1}\to A$ 的映射的类型，这些映射在初始顶点处为 $x$，在终端顶点处为 $y$。这种构造产生的类型不是范畴，因为这些态射的先验组合是不可能的。然而，如果限制自己研究 Rezk 类型，该理论允许我们证明关于 $(\infty,1)$-范畴的几个陈述。

除了类型不是范畴的问题之外，我们无法构造一个合理的范畴范畴。更糟糕的是，笔直化与反笔直化的 Grothendieck 构造在范畴意义上不起作用。我们将在第3.8节详细探讨这一点。

扩展 HoTT 的优点在于所有现有陈述仍然成立，这很好，因为 HoTT 已被证明是一个有用的理论。

\subsection{限制 $\Pi$-类型}

另一种方法是本质上限制 $\Pi$-类型的形成。这些是类型论中对应全称量词的概念，并推广了普通函数。并非所有依赖类型都应该被允许转化为 $\Pi$-类型。D.-C. Cisinski, B. Cnossen, K. Nguyen, and T. Walde [20] 以非类型论的方式探索了这一点，但应该可以从公理中提取出一个类型论。这种方法的优点在于它真正展示了综合范畴论。他们的工作甚至没有定义“类型”的概念。存在一个范畴的范畴，并且 Grothendieck 构造按预期工作。

这种方法的缺点是只有余笛卡尔纤维化支持幂对象（参见[20, 公理 N]）。这对应于类型论中 $\Pi$-类型的形成，因此 $\Pi$-类型仅对作为余笛卡尔族的依赖类型存在。这意味着特别地，并非所有 HoTT 中的陈述都可以转移，并且所有陈述都需要在依赖它们之前仔细检查（例如通过使用证明辅助工具证明它们）。关于这一理论的工作尚未完成，也没有证明辅助工具形式化这一理论，我们在此不再进一步研究它。

\subsection{圣杯}

由于两种方法都放弃了另一世界的一些东西，我们无法重现 HoTT 和 $\infty$-群胚的成功。范畴友好的限制 $\Pi$-类型的方法很可能产生更自然的范畴论环境。然而，形式化方法是否也能在那里表现良好仍有待观察。

\section{动机}

高阶范畴论中一个出现在许多地方的概念是余笛卡尔纤维化。它们在 D.-C. Cisinski, B. Cnossen, K. Nguyen, and T. Walde [20] 中扮演了重要角色，出现在三个公理中。在 E. Riehl 和 M. Shulman [2] 中，它们是第2.8节讨论的忠实 Grothendieck 构造的核心障碍的一部分。此外，余笛卡尔纤维化在高阶范畴论中通常具有重大意义，因为它们允许构造无穷范畴之间的函子，正如 M. Land [21] 所讨论的那样。

基于这些原因，我们在 \textsf{rzk} 中形式化了余笛卡尔纤维化，并将我们的工作上游集成到[3] 中。我们特别遵循了 U. Buchholtz 和 J. Weinberger [5] 以及 J. Weinberger [7] 的工作，他们将余笛卡尔纤维化作为 LARI 族的特例进行研究。

\section{综合 vs 分析}

我们已经提到，我们采用综合的方法做数学。它是分析方法的对偶——综合方法旨在公理化地产生一个行为正确的定义，而不是在集合论内部构造一个理论。一个著名的关于自然数的综合理论是皮亚诺公理。这些公理通过声明“0”存在来定义它；此外，对于任何自然数，都存在一个与其之前所有数不同的后继。从这些少数性质（以及归纳法等更多公理）出发，可以推导出许多关于自然数的有用陈述。相比之下，定义自然数的分析方法例如冯·诺依曼构造，它嵌入在 Zermelo-Fraenkel 集合论中。自然数被定义为集合，从 $0\coloneqq \emptyset$ 开始，$n$ 的后继为 $n\cup \{n\}$。加法、减法和任何其他运算都使用基本集合运算来定义。

分析方法允许使用现有理论，并表明数学概念存在于该理论中。然而，将它们嵌入另一个理论的需要使得它们本质上更加冗长。通常，人们必须“深入”定义，并使用关于具体表示的事实来得出所需陈述的证明。这在形式化数学时尤其如此。当在一个将实数构造为 Cauchy 序列等价类的理论中将 2 和 1 相加时，证明辅助工具有两种选择。要么它帮助解构序列，逐元素相加，然后观察到结果序列也是常数。要么，数学家必须执行这些额外的工作，而证明辅助工具只检查该工作的正确性。综合方法通过专门为概念构建理论来弥补这个问题，从而得到更清晰的证明和更大的普适性。在任何公理可以被分析证明的理论中，我们可以依赖从综合理论中证明的任何陈述。

\section{论文结构}

第2节介绍了同伦类型论，并涵盖了来自一般类型论的许多重要概念。我们遵循无公理基础计划[1] 和 R. Harper 等人[19] 的工作，并解释 \textsf{rzk} 如何实现 HoTT。第3节在此基础上，添加了 E. Riehl 和 M. Shulman [2] 的扩展类型，并介绍了来自范畴论的许多重要概念。在第4节中，我们讨论了正交纤维化和族，以便定义内族。这些是余笛卡尔纤维化所需的基本概念。我们还在本部分研究了理论中正方形的行为。接下来，我们在第5节、第6节和第7节中介绍我们论文的主要工作，其中我们首先介绍了始截面以及依赖始截面的概念。我们证明它们都等价于 LARI 伴随，U. Buchholtz 和 J. Weinberger [5] 使用 LARI 伴随来定义 LARI 族。在涵盖余笛卡尔纤维化之前的最后一步，我们研究了 LARI-胞腔，这是 J. Weinberger [7] 创造的一个概念，概括了余笛卡尔纤维化。我们证明这些等价于我们之前定义的 LARI-族。最后，在第7.4节中，我们收集了抽象劳动的成果，并讨论了 $i_0$-LARI 族作为余笛卡尔纤维化的具体实例化。最后，在第8节中，我们解释了如何使用 \textsf{rzk} 形式化我们的大部分工作。在我们的形式化过程中，我们发现了一个 \textsf{rzk} 中的 bug，维护者迅速修复了它。我们在附录 A 中包含了一个词汇表，其中包含用于查找符号记法的快速参考。

\chapter{同伦类型论}

在本节中，我们遵循 [1] 对同伦类型论（HoTT）进行总结。我们涵盖了在后续章节中将要用到的所有内容。我们还解释了某些概念在 rzk 中是如何实现的。关于具体例子、更多背景和更长的解释，我们请读者参阅 [1] 和 [19]。

本节非常注重形式正确性，我们经常明确提及一些在后面会忽略的细节。我们通过在证明辅助工具 rzk [4] 中证明了我们的陈述，来为后面部分的宽松处理提供理由。

\section{类型论基础}

HoTT 很大程度上基于 Martin-Löf [10] 发展的类型论。尽管从他最初的论文发表已经过去了几十年，并且许多其他数学家也做出了贡献，但最初的思想仍然在我们的理论中可见。对于本小节，我们主要遵循 R. Harper 等人 [19] 的讲座。

在开始写下我们的理论之前，我们需要涵盖几个概念。类型论起源于逻辑学，并与构造性数学和证明相关数学有联系。例如，我们不会假定排中律成立，因为它与我们后面要介绍的幺半等价性公理直接冲突。逻辑学与类型论之间存在一个基本的联系，由 Curry-Howard 同构描述。它将逻辑学和类型论中的概念一一对应起来：



该表是我们本节将介绍的各种类型的重要动机来源。它们代表了有用的数学理论应该具备的某些逻辑概念。

在深入定义这些类型的具体细节之前，我们首先解释如何建立形式语言，以便能够写出它们。

\subsection{推理规则与判断}

为了写下数学理论，存在几种不同的符号系统。例如，集合论使用一阶逻辑。由于类型论不建立在任何逻辑之上，我们使用逻辑学家通常用来描述一阶逻辑（以及许多其他类型的逻辑）的相同符号系统来引入其原始概念。

这个符号系统是一种形式语言，在其中我们定义判断和推理规则（通常简称为规则）。判断是“我们可以在理论内证明的陈述”。规则是我们从其他判断推导出判断的方式。一条规则具有以下语法：

\[
\frac{J_1 \quad J_2 \quad J_3 \quad \cdots \quad J_n}{J} \; R_1
\]

这条规则名为 $R_1$。$J_i$ 和 $J$ 是判断。横线上方的 $J_i$ 是规则的前提，底部的 $J$ 是结论。我们允许零前提的规则，这些规则本质上对应于我们理论的公理，因为它们没有前提条件。一条规则可用于在其所有前提成立的情况下推导出其结论。我们的判断中有变量，这些变量在规则中本质上是“全称量化的”。这不是理论内部量词，而是允许对这些规则进行句法替换的结果。本质上它们是宏，我们可以为变量代入任何值，并且它们对我们能代入的任何值都成立。

规则用于构建判断的证明。要证明一个判断，我们需要写下一系列步骤，其中只允许应用推理规则。最后我们应该得出所需判断成立的结论。逻辑学家为此发明了特殊的符号：证明树。一个简单的例子可以在示例2.1.1中找到。推理规则可以像乐高积木一样组合，将一条规则的结论插入另一条规则的前提。如果最终树的所有叶子都是横线（表示它们在没有前提的情况下证明了结论），我们就可以得出底部的判断可以从公理推导出来的结论。

\begin{example}
基本的证明树
\[
\frac{\frac{J_1}{J_1} \quad \frac{J_2}{J_2} \quad \frac{J_3}{J_3} \quad \frac{J_4}{J_4}}{\frac{J_5}{J_5}}
\]
\end{example}

对此的一个趣味性解读是像国际象棋这样的游戏。我们的公理是棋子的起始位置，规则是规定棋子如何移动的游戏规则。一个判断就是棋子的某个特定位置。一个证明是一系列规则，按顺序应用于起始位置后，得到我们想要证明的位置。

一个理论的所有规则都需要被证明为合理的，因为我们声明它们为真而不加证明。严格来说，我们还需要证明系统本身。关于这一点可以给出的论证超出了数学的范畴。我们将其留给哲学家，并满足于这样一个简单的事实：这个系统允许我们描述非常有趣和有用的数学概念。

\subsection{玩具逻辑}

在深入探讨我们类型论的具体规则和判断之前，我们先快速给出一个（不完整的）命题逻辑的例子。这是一个简单的判断和推理规则系统，将使读者有机会在我们本节后面更复杂的类型论规则之前，获得对逻辑语言更直观的理解。我们还将引入非常重要的蕴涵、上下文等概念。

我们首先需要描述玩具逻辑的判断是什么。我们拥有的第一种判断只是普通变量 $A,B,C,\ldots$；我们也允许带下标的变量。

下一种判断也会在我们后面的类型论中出现（尽管形式更复杂）。它被称为蕴涵，并且要求我们定义一个新的句法概念：上下文。一个上下文只是一个变量列表。我们使用大写希腊字母 $(\Gamma ,\Xi ,\Theta)$ 表示上下文。蕴涵则允许我们表示一个判断可以从一个变量列表推导出来。我们写作 $A_{1},\ldots ,A_{n}\vdash B$，其中 $A_i$ 是蕴涵的前提，$B$ 是结论。当我们写 $\Gamma \vdash B$ 时，读作“在上下文 $\Gamma$ 中，$B$ 成立”。蕴涵主要由以下两条规则支配：

\[
\frac{}{\Gamma\vdash A} \text{（自反性）} \qquad \frac{\Gamma\vdash A \quad \Gamma,A\vdash B}{\Gamma\vdash B} \text{（传递性）}
\]

第一条规则是自反性：我们逻辑中的所有陈述必须能从自身推出。第二条规则是传递性，它对应于将引理的证明代入假设该引理为真的证明中。除了这两条规则，还有三条规则允许操作蕴涵的上下文：

\[
\frac{\Xi\vdash B}{\Gamma,\Xi,\Theta\vdash B} \text{（ weakening ）} \qquad \frac{\Gamma,A,\Xi,A,\Theta\vdash B}{\Gamma,A,\Xi,\Theta\vdash B} \text{（ contraction ）} \qquad \frac{\Gamma\vdash A}{\pi(\Gamma)\vdash A} \text{（ permutation ）}
\]

第一条规则称为弱化，允许向上下文添加不必要的变量。第二条是收缩，允许丢弃重复的变量。最后一条规则是置换，允许重新排序上下文；其中 $\pi(\Gamma)$ 表示“将置换 $\pi$ 应用于 $\Gamma$ 的元素”。

我们以易于应用的方式表述了这些规则。使用以下不同的规则就足够了，但那样我们会有更多的中间步骤：

\[
\frac{\Gamma\vdash B}{\Gamma,A\vdash B} \text{（弱化'）} \qquad \frac{\Gamma,A,A\vdash B}{\Gamma,A\vdash B} \text{（收缩'）} \qquad \frac{\Gamma,A,B,\Xi\vdash C}{\Gamma,B,A,\Xi\vdash C} \text{（置换'）}
\]

这就完成了我们玩具逻辑中的所有判断，我们可以进入命题逻辑的第一个原始概念：合取。我们引入表达式的新句法。变量是表达式，给定两个表达式 $E_{1},E_{2}$，我们声明 $E_{1}\wedge E_{2}$ 也是一个表达式。为了赋予这个句法其语义含义（匹配合取的直觉），我们需要三条规则：

\[
\frac{\Gamma\vdash A \quad \Gamma\vdash B}{\Gamma\vdash A\wedge B} \wedge\text{I} \qquad \frac{\Gamma\vdash A\wedge B}{\Gamma\vdash A} \wedge\text{E}_1 \qquad \frac{\Gamma\vdash A\wedge B}{\Gamma\vdash B} \wedge\text{E}_2
\]

第一条规则称为引入规则，它允许我们在合取的两个部分都为真时构造一个合取。第二和第三条规则是消除规则，它们允许我们从合取中证明两个组成部分。

由于在所有这些规则中，上下文 $\Gamma$ 作为顶部和底部的唯一上下文出现，它不传递任何新信息。根据这个论点，我们可以省略它。因此，如果可能，我们也允许以局部规则形式书写规则：

\[
\frac{A \quad B}{A\wedge B} \wedge\text{I} \qquad \frac{A\wedge B}{A} \wedge\text{E}_1 \qquad \frac{A\wedge B}{B} \wedge\text{E}_2
\]

除了合取，我们的玩具逻辑中还有蕴含。给定两个表达式 $E_{1},E_{2}$，我们声明 $E_{1}\to E_{2}$ 也是一个表达式。引入和消除规则是：

\[
\frac{\Gamma,A\vdash B}{\Gamma\vdash A\Rightarrow B} \rightarrow\text{I} \qquad \frac{A\Rightarrow B \quad A}{B} \rightarrow\text{E}
\]

我们注意到引入规则必须用显式的上下文书写，因为我们使用了两个不同的上下文。消除规则通常称为“肯定前件”。

最后，我们允许带括号的表达式，所以如果 $E$ 是一个表达式，那么 $(E)$ 也是一个表达式。我们禁止任何导致运算顺序歧义的表达式（也可以定义运算之间的优先级以避免不必要的括号）。

使用这些规则，我们现在可以从空上下文证明陈述 $A\rightarrow (B\rightarrow (A\wedge B))$：

\[
\frac{\frac{}{A\vdash A} \text{W} \quad \frac{}{B\vdash B} \text{W}}{\frac{A,B\vdash A \quad A,B\vdash B}{A,B\vdash A\wedge B} \wedge\text{I}}
\]

除了这些概念，我们还需要允许书写定义。我们为此操作使用冒号和等号“$=$”。它为右边的内容引入了一个左边的缩写。例如：

\[
X \coloneqq A\rightarrow B
\]

这意味着以后任何时候看到 $X$，都应该用 $A\rightarrow B$ 替换它。替换时需要注意运算符优先级，我们在玩具逻辑中没有定义。

后面会再次出现的两个逻辑概念是析取、假和真。这两个常量有特别简单的规则：

\[
\frac{}{\top} \top\text{I}
\]

真引入说在任何上下文中我们都可以推导出真。它没有消除规则，我们不能从“真”证明任何东西。对偶地，假显然没有引入规则。它的消除规则假定假，并允许证明任何原始判断 $A$。

析取对读者来说应该很熟悉。引入规则也非常简单，反映了与合取的对偶性。一个潜在的新颖之处是消除规则。

\[
\frac{A}{A\vee B} \vee\text{I}_1 \qquad \frac{B}{A\vee B} \vee\text{I}_2 \qquad \frac{\Gamma,A\vdash C \quad \Gamma,B\vdash C}{\Gamma,A\vee B\vdash C} \vee\text{E}
\]

消除规则允许我们在给定 $A\vee B$ 的情况下，当我们能同时从 $A$ 和 $B$ 证明 $C$ 时，证明任何陈述 $C$。这里一个有趣的点是，这给我们的理论增加了一些模糊性，因为我们在调用这条规则时必须选择 $C$。当我们后面讨论乘积类型时，由于我们类型论中的证明项，这种模糊性会消失。在我们转向类型论蕴涵之前，我们给出一个关于类型论中推理规则含义的非常重要的评论。

\subsection{规则与类型检查}

类型论中的推理规则扮演的角色比逻辑学中重要得多，因为它们不直接用于构造证明。类型论中的规则决定哪些项是有效的；这称为类型检查。然而，这些规则连同项的句法被设置成可以自动应用的方式。

由于这个性质，类型论自然适合借助计算机进行形式化。类型检查可以通过一个程序实现，它让用户提供有效的项。

在本论文的其余部分，我们将不再显式地使用这些规则。相反，我们声称某个项是有效的，并且由于检查这一事实只是一个算法练习，我们在本节之后不再给出证明树。

这是类型论的优点，也是相对于逻辑学的一大进步。证明树是一个笨拙的对象——项则轻便得多，也更容易写下来（想出正确的项仍然很困难，就像想出一个证明很困难一样）。

这一点得到了进一步证实：在我们用来形式化结果的证明辅助工具 rzk [4] 中，没有办法显式地使用推理规则。它们都由类型检查器完全实现，对用户不可见。只有当用户形成了一个无法通过类型检查的项，需要报告错误时，类型检查才显示给用户。rzk 会打印一个步骤的堆栈跟踪，解释预期的情况。我们在第8节更详细地解释这一点。

\subsection{类型论蕴涵}

我们首先从一些句法开始定义类型论，然后引入上下文。它们比在命题逻辑中扮演着更重要的角色，因为它们包含更多信息：它们允许我们定义依赖类型和项，这在我们理论中起着至关重要的作用。

我们写 $\Gamma\ \mathsf{ctx}$ 来声明一个上下文。我们用 $\cdot$ 表示空上下文，当它的缺失会引起混淆时使用。为了描述什么是上下文，我们还需要类型声明的句法构造。我们写 $\Gamma \vdash A\ \Type$ 来表示 $A$ 在上下文 $\Gamma$ 中是一个类型。此外，我们用 $a:A$ 表示项声明，说明项 $a$ 具有类型 $A$。上下文的创建则由以下规则支配：

\[
\frac{}{\cdot\ \mathsf{ctx}} \quad \frac{\Gamma\ \mathsf{ctx} \quad \Gamma \vdash A\ \Type}{(\Gamma, x:A)\ \mathsf{ctx}}
\]

\begin{remark}
我们注意到，将形成上下文的规则指定为推理规则在技术上使我们的类型论成为一个多类理论。因为当我们后面写下一条推理规则时，那里的上下文形成需要遵循我们上面写的规则。这并不罕见，事实上 [19] 在没有提及的情况下就这样做了。
\end{remark}

当我们为某个判断 $J$ 写 $a:A\vdash J$ 时，这意味着变量 $a$ 被允许出现在该判断中。当我们后面添加宇宙的概念时，在使用 $J$ 时可以摆脱这种隐式依赖。

接下来我们给出蕴涵的自反性和传递性规则。类型论中的蕴涵传递性规则采用了略有不同的形式，需要一个新的句法概念，称为替换。它也更常被称为替换规则：

\[
\frac{}{\Gamma, x:A, \Xi \vdash x:A} \text{（自反性）} \qquad \frac{\Gamma \vdash a:A \quad \Gamma, x:A, \Xi \vdash J}{\Gamma, \Xi[x\mapsto a] \vdash J[x\mapsto a]} \text{（替换）}
\]

句法 $[x\mapsto a]$ 将替换应用于直接写在它前面的变量。替换将 $x$ 的所有出现改为 $a$。我们允许写 $M[x\mapsto a,y\mapsto b]$ 代替 $(M[x\mapsto a])[y\mapsto b]$。

\begin{remark}
替换由于需要区分自由变量和绑定变量而使理论复杂化。我们完全忽略这个技术细节，因为有一个规范的解决方案：不替换绑定变量。这可以由计算机实现。或者，我们可以只使用每个变量一次。为了方便，我们允许重复使用变量。要找到变量的类型，在向上遍历子表达式树时，只考虑最近的声明。
\end{remark}

如果一个变量不需要为了项通过类型检查而将其添加到上下文中，那么该变量是绑定的。自由变量是需要存在于上下文中才能使项通过类型检查的变量。

我们没有单一的弱化规则，而是给出一个规则模式。对于每个适合该句法位置的判断 $J$，我们都有一个采用以下形式的推理规则：

\[
\frac{\Gamma,\Xi \vdash J \quad \Gamma \vdash A\ \Type}{\Gamma, x:A, \Xi \vdash J} \text{（弱化）}
\]

收缩和置换也必须是规则模式：

\[
\frac{\Gamma \vdash A\ \Type \quad \Gamma, x:A, y:A, \Xi \vdash J}{\Gamma, z:A, \Xi[x\mapsto z,y\mapsto z] \vdash J[x\mapsto z,y\mapsto z]} \text{（收缩）}
\]

\[
\frac{\Gamma \vdash A\ \Type \quad \Gamma \vdash B\ \Type \quad \Gamma, a:A, b:B, \Xi \vdash J}{\Gamma, b:B, a:A, \Xi \vdash J} \text{（置换）}
\]

置换规则仅在 $a$ 不出现于 $B$ 且 $b$ 不出现于 $A$ 时才能应用。

\begin{remark}
与我们的玩具逻辑相比的新颖之处在于我们可以有依赖项/类型。例如，如果我们有一个自然数类型 $\mathbb{N}$，我们可以声明长度为 $n$ 的序列的依赖类型，其中 $n$ 是一个自然数。用我们的形式语言，我们可以通过判断来表达：
\[
A\ \Type, n:\mathbb{N} \vdash \mathrm{Seq}_n(A)\ \Type
\]
在 rzk 中，没有蕴涵。它使用 $\Pi$-类型和宇宙（我们稍后将介绍两者）来代替依赖类型。定义也允许指定一个上下文，该上下文会被翻译成 $\Pi$-类型。
\end{remark}

\begin{remark}
关于类型论的证明论观点是，类型的元素是类型所代表的陈述的证明。我们可以有同一陈述的多个不相等的证明，这使得我们的理论是证明相关的。注意，经典数学是证明无关的，从理论内部不可能区分同一陈述的两个不同证明。甚至一开始就不可能比较它们，因为一个判断要么成立，要么不成立。

一个拥有元素的类型被称为是居留的，它对应于经典数学中的一个真陈述。没有任何证明的类型被称为非居留的，它对应于假。
\end{remark}

我们玩具逻辑缺少的另一个概念是相等性，所以我们接下来介绍它。

\subsection{判断相等}

我们要介绍的下一个概念是判断相等。它也被称为定义相等。它描述了两个句法元素何时是“相同的”。这意味着在我们的整个理论中，我们应该能够在任何地方用一个替换另一个而不引起语义变化。我们需要将这种相等性与命题相等区分开来，命题相等我们稍后将使用普通的 $=$ 符号作为一种类型引入，而判断相等使用 $\equiv$。$=$-类型是同伦类型论中最重要的类型，它赋予了理论有趣的无限结构。

我们写 $\Gamma \equiv \Gamma'$ 表示两个上下文 $\Gamma$ 和 $\Gamma'$ 判断相等。对于类型，如果 $\Gamma \vdash A\ \Type$ 且 $\Gamma \vdash A'\ \Type$，我们写 $\Gamma \vdash A\equiv A'$ 表示 $A$ 和 $A'$ 在上下文 $\Gamma$ 中判断相等。最后，我们也允许对类型的元素进行相等比较，所以给定 $\Gamma \vdash a:A$ 和 $\Gamma \vdash a':A$，我们写 $\Gamma \vdash a\equiv a'$ 表示 $a$ 判断等于 $a'$。因此我们还需要 $\equiv$ 的三条自反性规则。

\[
\frac{\Gamma\ \mathsf{ctx}}{\Gamma\equiv\Gamma} \qquad \frac{\Gamma\vdash A\ \Type}{\Gamma\vdash A\equiv A} \qquad \frac{\Gamma\vdash a:A}{\Gamma\vdash a\equiv a}
\]

判断相等是一个非常基本的概念，因此它需要与我们添加到形式语言中的其他每个概念具有兼容性规则。对于类型（我们尚未介绍），这些规则称为 $\eta$ 和 $\beta$ 规则，我们在定义每种类型时给出它们。为了使判断相等与蕴涵兼容，我们需要以下规则（我们不给出名称，因为我们不显式使用它们）：

\[
\frac{\Gamma\vdash J \quad \Gamma\equiv\Gamma'}{\Gamma'\vdash J} \qquad \frac{\Gamma\vdash A\ \Type \quad \Gamma\vdash A'\ \Type \quad \Gamma\vdash A\equiv A'}{\Gamma\equiv\Gamma'} \qquad \frac{\Gamma\vdash a:A \quad \Gamma\vdash a':A \quad \Gamma\vdash a\equiv a' \quad \Gamma,x:A\vdash B\ \Type}{\Gamma\vdash B[x\mapsto a]\equiv B[x\mapsto a']}
\]

\[
\frac{\Gamma\vdash a:A \quad \Gamma\vdash a':A \quad \Gamma\vdash a\equiv a' \quad \Gamma,x:A\vdash b:B}{\Gamma\vdash b[x\mapsto a]\equiv b[x\mapsto a']}
\]

第一条规则是一个模式，允许我们从相等的上下文中证明相同的事情。第二条规则描述了我们如何证明上下文的相等性。第三和第四条规则表明，在替换的值相等的情况下，依赖类型和项是相等的。

在 rzk 中，判断相等也使用 $\equiv$ 符号。它的类型检查器中有相同的替换规则，因此允许在所有情况下交换判断相等的值。

\subsection{一般句法}

我们使用 $\coloneqq$ 进行定义，写 $X\coloneqq a$ 引入 $X$ 作为 $a$ 的简写符号（这里意义不大，但 $a$ 可以是任何表达式，而 $X$ 需要是一个标识符）。

在 rzk 中使用了非常相似的符号：
\begin{verbatim}
#def X : A := a
\end{verbatim}

这里除了 $X$ 的表示之外，还需要指定它的类型。在指定类型的 : A 前面，我们允许放置上下文变量。句法要求将它们放在括号中：
\begin{verbatim}
#def X (x : A) : B := a
\end{verbatim}

在我们的理论中，我们也有带括号的表达式，所以如果 $E$ 是一个表达式，$(E)$ 也是一个有效的表达式。在某些情况下，我们允许省略严格来说语法上需要的括号。例如，将函数应用于元组时，我们允许写 $f(a,b)$ 而不是 $f((a,b))$。这与将带有两个参数的函数应用于两个值有歧义，但无论如何这些函数类型通过柯里化是等价的，所以我们不会遇到问题。

\subsection{单位类型}

我们介绍的第一个类型将是非常简单的。它被称为单位类型，写作 $\mathbb{1}$，它恰好由一个元素组成，我们使用语法 $\ast_{\mathbb{1}}$。

\[
\frac{}{\mathbb{1}\ \Type} \qquad \frac{}{\ast_{\mathbb{1}}:\mathbb{1}}
\]

由于我们也会将 $\ast$ 的表示法用于其他只有一个成员的类型，如果有歧义或者我们想更明确地说明我们讨论的是哪个 $\ast$，我们会写 $\ast_{\mathbb{1}}$。

这个类型只有一个 $\eta$ 规则，没有 $\beta$ 规则：

\[
\frac{x:\mathbb{1}}{x\equiv\ast} \; \mathbb{1}\text{-}\eta
\]

这表明单位类型的任何值都必须判断等于 $\ast$。

这个类型不是特别有趣，因为它是最平凡的类型。它表示逻辑中 $\top$/真的值。从范畴论来看，它是 Cat 中的终端对象。只有一个函子映射到 $\mathbb{1}$（注意这个唯一性取决于判断相等！）。

在 rzk 中，它名为 Unit，该类型的唯一值称为 unit。它也有 $\eta$ 规则，所以每当我们有一个依赖于 x: Unit 的定义时，我们知道 x ≡ unit。

\subsection{空类型}

空类型表示假，写作 $\mathbb{0}$，没有元素。在它存在的情况下，我们可以证明任何类型的居留性。

\[
\frac{}{\mathbb{0}\ \Type} \qquad \frac{\Gamma \vdash a:\mathbb{0}}{\Gamma \vdash \mathrm{ind}_{\mathbb{0}}(a):A} \; \mathbb{0}\text{-E}
\]

没有 $\beta$ 或 $\eta$ 规则。

我们注意到这个类型在 rzk 中不可用。从范畴论来看，它是 Cat 中的始对象，由于消去规则，我们总是可以像预期那样定义到任何任意类型的函数。

\subsection{$\Sigma$ 类型}

我们介绍的 Martin-Löf 类型论中的下一个类型是 $\Sigma$-类型。它是对乘积类型的推广，我们首先介绍乘积类型。形成乘积类型的句法是 $A\times B$。我们还有两个投影函数，称为 $\pi_{1}$ 和 $\pi_{2}$。乘积类型的规则（以局部规则形式给出）是：

\[
\frac{A\ \Type \quad B\ \Type}{A \times B\ \Type} \times\text{-F} \qquad \frac{a:A \quad b:B}{(a,b):A \times B} \times\text{-I} \qquad \frac{p:A \times B}{\pi_1(p):A} \times\text{-E}_1 \qquad \frac{p:A \times B}{\pi_2(p):B} \times\text{-E}_2
\]

这些规则从左到右依次称为乘积类型的形成规则、引入规则和消除规则。大多数类型都有所有这三种规则。形成规则描述了我们什么时候甚至可以谈论乘积类型。我们从规则中省略了所有非依赖的“类型”判断，因为它们非常分散注意力且不包含新信息（总是可以从上下文中推断出来）。第二条规则是引入规则。它说明我们如何引入乘积的值。接下来的两条规则是其对应物，消除规则。它们规定了可以用乘积类型的值做什么。注意我们如何使用其局部规则形式指定所有这些规则，因为任何地方都没有使用上下文。

注意，虽然 $\pi_1$ 和 $\pi_2$ 看起来像函数，但严格来说它们是只是行为类似于函数的句法概念（我们尚未定义函数）。这种区别在实践中无关紧要，在本论文中我们就把它们当作函数对待。

用于与 $\equiv$ 兼容的 $\beta$ 和 $\eta$ 规则是：

\[
\frac{a:A \quad b:B}{\pi_1((a,b)) \equiv a} \times\beta_1 \qquad \frac{a:A \quad b:B}{\pi_2((a,b)) \equiv b} \times\beta_2 \qquad \frac{x:A \times B}{(\pi_1(x), \pi_2(x)) \equiv x} \times\eta
\]

另一种写这些规则的方式是通过等式，因为规则的前提从结论中显而易见（本质上我们只需要对所有变量进行全称量化，然后运行类型推断算法）：

\[
\pi_1((a,b)) \equiv a \qquad \pi_2((a,b)) \equiv b \qquad (\pi_1(x), \pi_2(x)) \equiv x
\]

乘积类型是类型论中对应合取的概念。这是因为当证明 $A \times B$ 的居留性时，需要给出 $A$ 的证明和 $B$ 的证明。

$\Sigma$-类型提供的推广是依赖性。第二个分量的类型可以依赖于第一个分量的值。例如，使用 $\Sigma$-类型，我们可以形成一个单一的类型，其中包含任意长度的序列。

\begin{definition}[$\Sigma$-类型]
$\Sigma$-类型的句法是 $\Sigma(x:A,B)$，其中 $A$ 是一个类型，$B$ 是在变量 $x:A$ 上的依赖类型。这个构造绑定变量 $x$，该变量允许在 $B$ 中出现。以后当我们有宇宙时，这种依赖性会更清晰地标记为 $\Sigma(x:A,B(x))$，其中 $B$ 是从 $A$ 到宇宙的函数，写作 $B:A \to \mathcal{U}$。

我们有下面的形成、引入和消除规则：

\[
\frac{\Gamma\vdash A\ \Type \quad \Gamma,x:A \vdash B\ \Type}{\Gamma\vdash \Sigma(x:A,B)\ \Type} \qquad \frac{\Gamma,x:A \vdash B\ \Type \quad \Gamma\vdash a:A \quad \Gamma\vdash b:B[x\mapsto a]}{\Gamma\vdash (a,b):\Sigma(x:A,B)}
\]

\[
\frac{\Gamma\vdash \Sigma(x:A,B)\ \Type \quad \Gamma\vdash p:\Sigma(x:A,B)}{\Gamma\vdash \pi_1(p):A} \qquad \frac{\Gamma\vdash \Sigma(x:A,B)\ \Type \quad \Gamma\vdash p:\Sigma(x:A,B)}{\Gamma\vdash \pi_2(p):B[\pi_1(p)/x]}
\]

此外，$\beta$ 和 $\eta$ 规则与乘积类型相同（但现在如果我们进行类型检查，有 $b:B[x\mapsto a]$ 而不是 $b:B$）：

\[
\pi_1((a,b))\equiv a \qquad \pi_2((a,b))\equiv b \qquad (\pi_1(x),\pi_2(x))\equiv x
\]

我们允许写 $\Sigma(a:A,b:B,C)$，它应解释为 $\Sigma(a:A,\Sigma(b:B,C))$。对于项我们也做同样处理，所以 $(a,b,c)\equiv (a,(b,c))$。在内部类型占用大量垂直空间的情况下，我们也使用更大的“$\sum$”符号。

类型论文献中经常使用略有不同的 $\Sigma$-类型符号：
\[
\sum_{x:A}B \equiv \Sigma(x:A,B)
\]

我们选择自己的符号，因为它避免了索引漂移，例如比较下面两种表示法的类型：
\[
\sum_{x:\sum_{y:\sum_{a:A}B}C}D \quad \text{vs} \quad \Sigma(x:\Sigma(y:\Sigma(a:A,B),C),D)
\]
\end{definition}

\begin{example}
我们给出一个 $\Sigma$-类型的基本例子。假设我们有自然数类型 $\mathbb{N}$ 和长度为 $n:\mathbb{N}$ 的序列类型 $\mathrm{Seq}_n(A)$，其中 $A$ 是一个类型。

我们现在可以形成 $\Sigma$-类型：
\[
\Sigma(n:\mathbb{N},\mathrm{Seq}_n(A))
\]
这个类型是任意长度序列的类型。它的一个元素由长度 $n:\mathbb{N}$ 和一个该长度的序列组成。

这个类型在逻辑上等价于陈述“存在一个自然数，使得存在一个 $A^n$ 中该长度的序列”。这个陈述非常平凡，表明该类型的所有有趣性质必须来自元素之间的差异。

我们的符号对和类型的两个分量赋予相同的优先级，而求和符号使得第二个分量重要得多。此外，我们的符号更接近 rzk 中使用的符号：
\begin{verbatim}
\Sigma (a : A), B
\end{verbatim}

这种符号的一个缺点是，在嵌套层次较高时，很难辨别哪个括号属于哪个 $\Sigma$。我们使用更大的括号、向量表示法或特殊换行/格式来防止影响可读性。在编写 rzk 代码时，[3] 使用的特殊格式确保任何嵌套层次都易于识别。
\end{example}

正如我们之前提到的，$\Sigma$-类型对应于逻辑中的存在量词：为了给出 $\Sigma(x:A,B)$ 的一个元素（从而证明 $\exists x:A,B$ 成立），首先必须提供一个见证 $a:A$，然后是一个证明 $b:B[x\mapsto a]$。这个证明 $b$ 表明对于我们的见证 $a$，类型 $B$（其中 $x$ 被替换为 $a$）是居留的（逻辑上说：见证 $a$ 满足条件 $B$）。

为了恢复普通的独立乘积类型，我们只需要为 $B$ 选择一个（在 $x$ 上）非依赖的类型。在这种情况下，我们允许写 $A\times B\equiv \Sigma(x:A,B)$（其中 $x$ 不在 $B$ 中自由出现）。

$\Sigma$-类型的范畴论版本称为“依赖和”[22]。它不太常见，因为它直接使用我们后面称为类型族的纤维化版本来工作。这使得它使用时更加冗长。它推广了余积的概念。在我们的类型论中，我们可以使用 $\Sigma$ 类型形成任何范畴极限。

\begin{remark}
有时会遇到非常大的 $\Sigma$-类型。在这种情况下，将元素以及类型本身写作向量可以显著提高可读性，特别是当每个分量的项包含括号和许多其他子表达式时。向量总是从上到下列出项目。例如：
\[
\Sigma(a:A,\Sigma(b:B,\Sigma(c:C,D))) \quad \text{vs} \quad \Sigma \begin{pmatrix} a:A \\ b:B \\ c:C \\ D \end{pmatrix}
\]
我们另外还允许在 $\Sigma$-类型中省略元素的名称，当它们不在后面的类型中出现时。这也可以使用 $\times$ 来写，但在某些情况下更节省空间。
\end{remark}

\subsection{$\Pi$ 类型}

接下来我们定义 $\Pi$-类型。这是函数类型的推广，也可以解释为乘积类型，因此得名。它们在逻辑上的对应物是全称量词 $\forall x.P$；如果有一个类型 $A$ 的值 $x$，那么 $\forall x.P$ 的证明“允许我们获得 $P[x\mapsto a]$ 的证明”。符号与 $\Sigma$-类型非常相似，但使用大写的 $\Pi$：$\Pi(a:A,B)$。该类型的值使用经典的 lambda 表示法 $\lambda a:A.b$，其中 $x$ 是类型 $A$ 的变量，点在之后绑定，$b$ 是一个项。

乘积解释是：$\Pi(i:I,E_i)$ 的一个值由每个 $i:I$ 的 $e_i:E_i$ 组成。对我们来说，$\Pi$-类型作为依赖函数类型的类型论解释是最有用的。这是因为它们表现出的行为类似于连续映射——只允许它们的值以一致的方式变化。我们将在本节后面定义可缩性概念时，以及在第五节的更后面，再次观察到这种有些模糊的行为。

\begin{definition}[$\Pi$-类型]
\[
\frac{\Gamma\vdash A\ \Type \quad \Gamma,x:A \vdash B\ \Type}{\Gamma\vdash (x:A)\to B\ \Type} \qquad \frac{\Gamma,x:A \vdash b:B}{\Gamma\vdash \lambda x.b:(x:A)\to B}
\]
\[
\frac{\Gamma\vdash f:(x:A)\to B \quad \Gamma\vdash a:A}{\Gamma\vdash f(a):B[x\mapsto a]}
\]

$\beta$ 和 $\eta$ 规则是：
\[
(\lambda x.b)(a)\equiv b[x\mapsto a] \qquad (\lambda x.b(x))\equiv b
\]
$\eta$ 规则（第二个）仅在 $x$ 不在 $b$ 中自由出现时成立。

我们允许写 $(a:A,b:B,c:C)\to D$ 代替 $(a:A)\to (b:B)\to (c:C)\to D$。在独立情况下，我们也写 $A\to B\to C\to D$，这不应与从 $A$ 到 $D$ 的复合函数的通常数学含义混淆。当我们要写复合时，我们总是在箭头上方放置名称。这种 $\Pi$-类型的符号再次类似于 rzk 中的符号。传统上，人们会为此类型使用乘积符号：
\[
\prod_{a:A}B
\]
然而，这也有索引漂移的问题，因此我们更喜欢我们的符号。

对于 lambda 表达式，我们也允许在变量后放置类型以帮助读者理解独立 lambda 表达式的类型：$\lambda x:A.b$。我们也允许对 $\Sigma$-类型的参数进行模式匹配。假设我们想写一个域为 $\Sigma(a:A,B)$ 的类型，我们这样写：$\lambda (a,b).x$，其中 $x$ 可以是任何项。

与 $\Sigma$-类型一样，在写类型和定义映射时允许类似向量的表示法：
\[
\Pi \begin{pmatrix} a:A \\ b:B \\ C \end{pmatrix} \rightarrow D \qquad \lambda \begin{pmatrix} a:A \\ b:B \end{pmatrix} . x
\]

类似于 $\Sigma$-类型，我们可以恢复标准非依赖函数类型的概念（我们直接写作 $A \to B$）。

范畴论上，这是依赖积[23]，类似于依赖和，很少被教授，并且是用纤维化而非族来表述的。
\end{definition}

\begin{definition}[恒等函数]
对于任何类型 $A$，我们写作 $\operatorname{id}_A$ 表示恒等函数 $A \to A$，定义为 $\lambda a.a$。

在 rzk 中，它是这样定义的：
\begin{verbatim}
#def id (A : U) (a : A) : A := a
\end{verbatim}
\end{definition}

rzk 中 $\Pi$ 类型的句法是 $(a : A) \to B$，对于 lambda 表达式，它使用反斜杠后跟参数，然后是一个箭头，再是 lambda 的主体：$\\a \to a$。

A : U 的部分声明 A 属于类型 U，即类型的宇宙。接下来我们将引入它作为“A 类型”判断的替代。在下一节之后，我们将能够定义一个单一的恒等函数，而不需要为每个类型都给出模式定义。

\begin{definition}[函数复合]
给定两个可复合的函数 $f:A\to B, g:B\to C$，我们定义：
\[
g\circ f \coloneqq \lambda a.g(f(a))
\]
\end{definition}

\end{document}